\chapter{Lyapunov 稳定性理论：从直接法到 ISS 与反步法}
\label{ch:lyapunov-stability}

% ════════════════════════════════════════════════════════════════
%  控制部 第 D 章（归卷四《控制理论基础》）。
%  源：Robotics_Tutorial/01_数学/40_控制理论/70_Lyapunov稳定性理论.md（CC-BY）。
%  定位：ch:control（导论，会用）→ ch:lyapunov-basics（C0，做透基础）→ 本章（研究生密度严格化）。
%  右扰动/Hamilton 约定与全书一致；本章为欧氏状态空间 R^n 上的稳定性理论，几乎不涉流形扰动。
% ════════════════════════════════════════════════════════════════

\begin{practice}[前置自测]
开始之前，先试答下列问题。若已能流畅作答，可只速读本章的定理速查表与小结；若卡住，正是本章要为你打磨的。
\begin{enumerate}
  \item 对线性系统 $\dot{\mathbf{x}}=\mathbf{A}\mathbf{x}$，特征值的实部如何决定稳定性？非线性系统 $\dot{\mathbf{x}}=f(\mathbf{x})$ 为何不能照搬这一判据？
  \item “收敛到原点”与“渐近稳定”是同一回事吗？若不是，缺了哪个条件？
  \item 一个有阻尼的单摆，总能量沿轨迹只减不增，但在最高点（速度为零）耗散为零——只凭 $\dot V\le 0$ 半负定，你能断定它最终停在竖直向下吗？需要什么额外工具？
  \item 真实机器人总有扰动（电机噪声、估计误差）。此时不能要求状态趋于零，那“鲁棒稳定”该如何定量刻画？
  \item 若你怎么也找不到一个满足条件的 $V$，能据此断言系统不稳定吗？
\end{enumerate}
\end{practice}

\section*{本章目标}
学完本章，你应当能够：
\begin{enumerate}
  \item \textbf{用比较函数} $\mathcal{K}/\mathcal{K}_\infty/\mathcal{KL}$ 统一刻画三级稳定性，消除繁琐的 $\varepsilon$-$\delta$ 摆弄；
  \item \textbf{完整证明} Lyapunov 直接法与 LaSalle 不变原理（含 $\omega$-极限集的不变性论证），而不止于会用；
  \item \textbf{把指数稳定与 Lyapunov 方程焊接}，理解它是 Riccati 的退化、是 LQR 闭环值函数的来源；
  \item \textbf{系统构造} Lyapunov 函数（能量法、二次型、Krasovskii、变量梯度、SOS、Neural），并知道每种方法的边界与陷阱；
  \item \textbf{陈述逆定理}（Massera/Kurzweil/Lin–Sontag–Wang），理解“稳定 $\Leftrightarrow$ 存在 $V$”的哲学完备性；
  \item \textbf{用 ISS 框架}分析带扰动的非线性系统鲁棒性，并用小增益定理判断互联/级联稳定性；
  \item \textbf{用反步法}递归设计镇定控制律，并完整证明机器人 PD+重力补偿的渐近稳定性；
  \item \textbf{估计吸引域}（sublevel set / SOS / Zubov 方程），并处理非自治系统（UAS + Barbalat 引理）。
\end{enumerate}

\subsection*{与控制导论、稳定性基础两章的分工}
本章是控制部稳定性主线的\emph{第三层}。第一层是 \cref{ch:control}（控制导论）：它给出三级稳定性定义（\cref{def:ctrl-lyap-stable}）、直接法（\cref{thm:ctrl-lyap-direct}）、LaSalle（\cref{thm:ctrl-lasalle}）、Lyapunov 方程（\cref{thm:ctrl-lyap-eq}）、间接法与 Barbalat 一句话提及（\cref{thm:ctrl-indirect}）的\emph{陈述与用法}——“会用”。第二层是 \cref{ch:lyapunov-basics}（稳定性基础）：它把定号性、间接法、Lyapunov 方程的计算、Krasovskii 法与变量梯度法\emph{做透到本科密度}。本章承接这两层，定位是\textbf{研究生密度的严格化与大幅扩展}：给直接法/LaSalle 补\emph{完整证明}与每条件的必要性反例；把指数稳定与 Riccati、逆定理的存在性、ISS/小增益/级联/反步这套\emph{鲁棒与设计}语言系统建立起来。凡前两章已推导透的结论，本章一律交叉引用指回、不重复推导，只补本章视角的深化。一句话：导论“会用”，基础“会算”，本章“会证、会扩”。

\begin{note}
\textbf{本章记号与约定.}\;
状态 $\mathbf{x}\in\mathbb{R}^n$、向量场 $\mathbf{f}(\mathbf{x})$ 用粗体（与 \cref{ch:control} 一致）。自治系统 $\dot{\mathbf{x}}=\mathbf{f}(\mathbf{x})$，$\mathbf{f}(\mathbf{0})=\mathbf{0}$，原点为平衡点。$V$ 沿轨迹的导数（Lie 导数）记 $\dot V=\nabla V\cdot\mathbf{f}=\sum_i\frac{\partial V}{\partial x_i}f_i$。比较函数类 $\mathcal{K},\mathcal{K}_\infty,\mathcal{KL}$ 在 \cref{sec:lyap-comparison} 定义后全章及后续 ISS/CLF/MPC 章复用。机械系统沿用 $\mathbf{M}(\mathbf{q})\ddot{\mathbf{q}}+\mathbf{C}(\mathbf{q},\dot{\mathbf{q}})\dot{\mathbf{q}}+\mathbf{g}(\mathbf{q})=\boldsymbol{\tau}$，增益矩阵统一记 $\mathbf{K}_p,\mathbf{K}_d\succ0$。本章是欧氏空间稳定性理论，几乎不涉 $\SO(3)/\SE(3)$ 流形扰动，故无右扰动记号冲突。
\end{note}

% ════════════════════════════════════════════════════════════════
\section{为什么 Lyapunov 是非线性稳定性的根基}
\label{sec:lyap-intro}

\paragraph{动机：一个无法回避的问题。}
设你为四足机器人设计了控制器，仿真中完美行走。但你如何\emph{证明}闭环系统不会突然发散？如何证明踩到一块石头（一个微小扰动）不会让它摔倒？这就是稳定性分析的核心。对线性系统 $\dot{\mathbf{x}}=\mathbf{A}\mathbf{x}$ 答案干净：看 $\mathbf{A}$ 的特征值是否都在左半平面（\cref{ch:lyapunov-basics} 的间接法）。但对 $\dot{\mathbf{x}}=\mathbf{f}(\mathbf{x})$，“特征值”不再适用——系统行为随状态而变，不同区域可能有完全不同的定性特征。

\paragraph{反面：没有 Lyapunov 方法会怎样。}
在 1892 年之前，分析非线性稳定性的唯一办法是\emph{显式求解}微分方程再验证解趋于零。但绝大多数非线性 ODE 没有闭式解：单摆 $\ddot\theta+\sin\theta=0$ 需椭圆函数，任何含摩擦、柔性、接触的机器人动力学都不可能解析求解。你永远“算不出”一个复杂系统的轨迹去验证收敛。

\begin{insight}
Lyapunov 方法的革命性在于\textbf{不需要解微分方程}：它把“求解”问题换成“构造”问题——找到一个广义能量 $V(\mathbf{x})$ 满足适当条件即可。这就像不需要知道水球在碗中的精确轨迹，只要知道“海拔在降低”，就足以断定它最终到达碗底。对非线性系统，这往往意味着从“不可能分析”变为“可以分析”。
\end{insight}

\paragraph{历史：三次跃迁。}
Lyapunov 1892 年博士论文《运动稳定性的一般问题》\cite{lyapunov1992general} 提出两种方法：第一方法（间接法）靠线性化判断，第二方法（直接法）靠能量函数直接判断——后者才是真正的突破。此后的发展可概括为三次跃迁：1960 年 LaSalle 不变集原理\cite{lasalle1960extensions} 把判据从 $\dot V<0$ 放宽到 $\dot V\le0$（救命）；1989 年 Sontag 的 ISS 理论\cite{sontag1989smooth} 把稳定性升级为带扰动的鲁棒性语言（统一）；2000 年后 Parrilo 的 SOS\cite{parrilo2000structured} 与 Neural Lyapunov\cite{chang2019neural} 把 $V$ 的构造\emph{计算化}（自动化）。本章的地图正沿这条线索铺开：先以比较函数立语言（\cref{sec:lyap-comparison}），证透直接法与 LaSalle（\cref{sec:lyap-direct,sec:lyap-lasalle}），打通指数稳定、Chetaev、构造法、逆定理（\cref{sec:lyap-exp,sec:lyap-chetaev,sec:lyap-construct-d,sec:lyap-converse}），再进入 ISS/小增益/级联/反步这套鲁棒与设计工具（\cref{sec:lyap-iss,sec:lyap-cascade,sec:lyap-backstepping}），落点于机器人 PD（\cref{sec:lyap-robot-pd}），最后处理吸引域与非自治系统（\cref{sec:lyap-roa,sec:lyap-nonauto}）。

% ════════════════════════════════════════════════════════════════
\section{比较函数体系与三级稳定性的等价刻画}
\label{sec:lyap-comparison}

\paragraph{动机：把稳定性编码为函数不等式。}
经典 $\varepsilon$-$\delta$ 定义虽严谨，但每证一个新性质都要重新摆弄 $\varepsilon,\delta$ 的关系。Hahn 1967\cite{hahn1967stability} 引入比较函数类，目的是\emph{把稳定性统一编码为函数不等式}，让不同层级只需改动一两项即可区分。这就像编程的类型系统：不必每次从位级别操作数据，而用抽象类型。比较函数是稳定性理论的“类型系统”，也是后续 ISS（\cref{sec:lyap-iss}）、CLF/CBF（\cref{ch:clf-cbf}）、MPC（\cref{ch:mpc-advanced}）的通用货币。

\begin{definition}[比较函数类 $\mathcal{K},\mathcal{K}_\infty,\mathcal{KL}$]
\label{def:lyap-class-k}
\begin{itemize}
  \item $\alpha:[0,a)\to[0,\infty)$ 属 $\mathcal{K}$（class-$\mathcal{K}$），若连续、$\alpha(0)=0$、严格递增；
  \item $\alpha:[0,\infty)\to[0,\infty)$ 属 $\mathcal{K}_\infty$，若 $\alpha\in\mathcal{K}$、定义域为 $[0,\infty)$ 且 $\lim_{r\to\infty}\alpha(r)=\infty$（无界）；
  \item $\beta:[0,\infty)\times[0,\infty)\to[0,\infty)$ 属 $\mathcal{KL}$，若对固定 $s$ 有 $\beta(\cdot,s)\in\mathcal{K}$、对固定 $r>0$ 关于 $s$ 严格递减且 $\lim_{s\to\infty}\beta(r,s)=0$。
\end{itemize}
\end{definition}

\noindent 直觉：$\mathcal{K}$ 是“从原点严格上升的曲线”（如 $r,\ r^2,\ \sqrt{r}$）；$\mathcal{K}_\infty$ 额外要求一路升到无穷（$\frac{r}{1+r}$ 属 $\mathcal{K}$ 但因有上界 $1$ 不属 $\mathcal{K}_\infty$）；$\mathcal{KL}$ 同时编码“初值越大响应越大”（$\mathcal{K}$ 性）与“随时间衰减到零”（L 性），典型为 $\beta(r,s)=r\,e^{-s}$。代数性质（证明留作练习）：$\mathcal{K}$ 对复合、求和封闭；$\alpha\in\mathcal{K}_\infty\Rightarrow\alpha^{-1}\in\mathcal{K}_\infty$（$\mathcal{K}$ 一般无逆）。

\begin{theorem}[三级稳定性的比较函数刻画]
\label{thm:lyap-comparison-char}
对 $\dot{\mathbf{x}}=\mathbf{f}(\mathbf{x})$、$\mathbf{f}(\mathbf{0})=\mathbf{0}$，原点
\begin{itemize}
  \item \textbf{Lyapunov 稳定} $\iff$ $\exists\alpha\in\mathcal{K}$（局部）使 $\|\mathbf{x}(t)\|\le\alpha(\|\mathbf{x}(0)\|),\ \forall t\ge0$；
  \item \textbf{渐近稳定} $\iff$ $\exists\beta\in\mathcal{KL}$ 使 $\|\mathbf{x}(t)\|\le\beta(\|\mathbf{x}(0)\|,t)$；全局即对所有 $\mathbf{x}(0)\in\mathbb{R}^n$ 成立；
  \item \textbf{指数稳定} $\iff$ $\exists k,\lambda>0$ 使 $\|\mathbf{x}(t)\|\le k\|\mathbf{x}(0)\|e^{-\lambda t}$（即 $\beta(r,s)=kr\,e^{-\lambda s}$ 的特例）。
\end{itemize}
\end{theorem}

\begin{insight}
$\mathcal{KL}$ 函数把“关于初值单调”与“随时间衰减到零”两个物理直觉\emph{打包成一个对象}。一旦掌握，所有进阶概念（ISS、iISS、ISpS、实际稳定）都是在这套语言上的微调——例如 \cref{sec:lyap-iss} 的 ISS 不等式只是在 $\beta(\|\mathbf{x}_0\|,t)$ 后多加一个 $\gamma(\|\mathbf{u}\|_\infty)$ 项。这正是把 $\varepsilon$-$\delta$ 一次性“封装”带来的复利。
\end{insight}

\begin{pitfall}[三个高频误区]
\textbf{(i) $\mathcal{K}$ 必须无界？}\;否。$\alpha(r)=1-e^{-r}$ 严格递增、$\alpha(0)=0$、连续，是 $\mathcal{K}$，但有上界 $1$，\emph{不}属 $\mathcal{K}_\infty$。只有 $\mathcal{K}_\infty$ 的下夹才能保证 $V$ 径向无界（\cref{sec:lyap-direct}）。
\quad\textbf{(ii) 收敛 $=$ 渐近稳定？}\;否。$\lim_{t\to\infty}\mathbf{x}(t)=\mathbf{0}$ 不蕴含渐近稳定——必须\emph{同时}有 Lyapunov 稳定。存在系统使所有轨迹最终趋零，却对任意小 $\varepsilon$ 都有任意小初值令轨迹先“暴冲”出 $\varepsilon$ 再回来；机器人先摔倒再站起来不算“稳定控制”。
\quad\textbf{(iii) 指数稳定最强？}\;否。有的系统\emph{有限时间}收敛，如 $\dot x=-\operatorname{sgn}(x)|x|^{1/2}$（指数 $\tfrac12\in(0,1)$），近原点比任何指数都快\cite{bhat2000finitetime}。指数稳定只是渐近稳定里一个特别干净的子类。又 GAS $\not\Rightarrow$ ES：$\dot x=-x^3$ 是 GAS（$V=x^2/2,\ \dot V=-x^4<0$），但解 $x(t)=x_0/\sqrt{1+2x_0^2t}$ 仅以 $O(1/\sqrt t)$ 多项式衰减。
\end{pitfall}

% ════════════════════════════════════════════════════════════════
\section{Lyapunov 直接法：完整证明与必要性}
\label{sec:lyap-direct}

\paragraph{承上启下。}
\cref{ch:control} 已陈述直接法（\cref{thm:ctrl-lyap-direct}）并给出“能量下降即稳定”的直觉，\cref{ch:lyapunov-basics} 已练熟其用法。本节补上导论\emph{略去的完整证明}——它揭示了为何“正定 $V$ + 半负定 $\dot V$”就锁死了稳定性，以及每个条件删去后究竟在哪一步崩溃。这套 $\varepsilon$-$\delta$ 与 $\omega$-极限集论证是 LaSalle（\cref{sec:lyap-lasalle}）、逆定理（\cref{sec:lyap-converse}）一切证明的母本。

\begin{theorem}[Lyapunov 直接法，完整版；深化 \cref{thm:ctrl-lyap-direct}]
\label{thm:lyap-direct-full}
设 $\mathcal{D}\subset\mathbb{R}^n$ 为含原点开集，$\mathbf{f}:\mathcal{D}\to\mathbb{R}^n$ 局部 Lipschitz，$\mathbf{f}(\mathbf{0})=\mathbf{0}$，$V\in C^1(\mathcal{D})$ 满足
\textbf{(a)} $V(\mathbf{0})=0,\ V(\mathbf{x})>0\ (\mathbf{x}\neq\mathbf{0})$；\textbf{(b)} $\dot V(\mathbf{x})\le0$ 于 $\mathcal{D}$。
则原点 Lyapunov 稳定。若进一步 \textbf{(c)} $\dot V(\mathbf{x})<0\ (\mathbf{x}\neq\mathbf{0})$，则原点渐近稳定。
\end{theorem}

\begin{proof}
\emph{稳定性.}\;给定 $\varepsilon>0$（小到闭球 $\bar B_\varepsilon\subset\mathcal{D}$）。球面 $S_\varepsilon=\{\|\mathbf{x}\|=\varepsilon\}$ 紧，$V$ 在其上正值连续，由 Weierstrass 定理 $m:=\min_{\|\mathbf{x}\|=\varepsilon}V(\mathbf{x})>0$（正定故 $m>0$）。由 $V(\mathbf{0})=0$ 与连续性，存在 $\delta\in(0,\varepsilon]$ 使 $\|\mathbf{x}\|<\delta\Rightarrow V(\mathbf{x})<m$。设 $\|\mathbf{x}(0)\|<\delta$，则 $V(\mathbf{x}(0))<m$；由 (b) $\dot V\le0$ 知 $V(\mathbf{x}(t))\le V(\mathbf{x}(0))<m,\ \forall t\ge0$。但 $V\ge m$ 于 $S_\varepsilon$，故轨迹\emph{永不触及} $S_\varepsilon$，即 $\|\mathbf{x}(t)\|<\varepsilon$ 恒成立——正是 Lyapunov 稳定的 $\varepsilon$-$\delta$ 定义。沿途用 sublevel set $\Omega_c=\{V\le c\}$（$c$ 小使其紧且 $\subset\mathcal{D}$）的正不变性：$V(\mathbf{x}(t))\le V(\mathbf{x}(0))\le c\Rightarrow\mathbf{x}(t)\in\Omega_c$。\\[2pt]
\emph{渐近稳定.}\;在稳定性基础上设 (c)。沿轨迹只要 $\mathbf{x}(t)\neq\mathbf{0}$ 即 $\frac{d}{dt}V(\mathbf{x}(t))=\dot V<0$，故 $V(\mathbf{x}(t))$ 严格递减且有下界 $0$，存在极限 $c_0\ge0$。反设 $c_0>0$：轨迹困于紧集 $\Omega_{V(\mathbf{x}(0))}$，其 $\omega$-极限集 $\omega(\mathbf{x}_0)$ 非空、紧、\emph{不变}。取 $\mathbf{p}\in\omega(\mathbf{x}_0)$，则 $V(\mathbf{p})=\lim_n V(\mathbf{x}(t_n))=c_0>0$，故 $\mathbf{p}\neq\mathbf{0}$，由 (c) $\dot V(\mathbf{p})<0$。但 $V$ 在不变集 $\omega(\mathbf{x}_0)$ 上恒为 $c_0$，从 $\mathbf{p}$ 出发的轨迹仍在其中，$V$ 应继续严格下降——矛盾。故 $c_0=0$。最后由正定性的比较函数下夹 $\alpha_1(\|\mathbf{x}\|)\le V(\mathbf{x})$（$\alpha_1\in\mathcal{K}$）得 $\alpha_1(\|\mathbf{x}(t)\|)\le V(\mathbf{x}(t))\to0$，即 $\mathbf{x}(t)\to\mathbf{0}$。
\end{proof}

\begin{table}[H]\centering\small
\caption{直接法各条件的必要性：删去则在何处崩溃}
\label{tab:lyap-direct-necessity}
\begin{tabular}{lll}
\toprule
条件 & 删去后崩在 & 反例 \\
\midrule
$V$ 正定 & 无法由 $V<m$ 推 $\|\mathbf{x}\|<\varepsilon$ & $V=x_1^2$（$x_2$ 方向无约束） \\
$\dot V\le0$ & 轨迹可逃出 $\Omega_c$ & 任一不稳定系统 $\dot x=x$ \\
$\dot V<0$（渐近） & 仅得稳定、轨迹绕而不收 & 简谐振子 $\dot x_1=x_2,\dot x_2=-x_1,\ \dot V=0$ \\
径向无界（全局） & 远处轨迹可沿“能量谷”逃逸 & $V=x^2/(1+x^2)$（上界 $1$） \\
\bottomrule
\end{tabular}
\end{table}

\begin{theorem}[全局渐近稳定；深化 \cref{thm:ctrl-lyap-direct}]
\label{thm:lyap-gas}
若 $V:\mathbb{R}^n\to\mathbb{R}$ 满足 \cref{thm:lyap-direct-full} 的 (a)(c) 且\emph{径向无界}（$\|\mathbf{x}\|\to\infty\Rightarrow V\to\infty$，等价于 $\exists\alpha_1\in\mathcal{K}_\infty,\ V(\mathbf{x})\ge\alpha_1(\|\mathbf{x}\|)$），则原点全局渐近稳定。
\end{theorem}

\begin{insight}
径向无界用在唯一一步：保证 $\Omega_c=\{V\le c\}$ 对\emph{所有} $c$ 都紧，于是无论初值多远，轨迹都被“框住”，渐近论证才能套用。删了它，$V$ 可能在远处“躺平”（如 $V=x^2/(1+x^2)<1$），$\Omega_1=\mathbb{R}^n$ 不紧，从远处出发的轨迹可能永不进入 $\dot V<0$ 起作用的内层——故只剩\emph{局部} AS。这与 \cref{ch:control} 的“承重假设三问”是同一回事，本章只是把“崩在哪步”落到了 $\Omega_c$ 的紧性上。
\end{insight}

\begin{pitfall}[计算 $\dot V$ 的两个常见错]
\textbf{(i) 忘了链式法则.}\;$\dot V$ 不是“对 $t$ 直接求导”，而是 $\nabla V\cdot\mathbf{f}=\sum_i\frac{\partial V}{\partial x_i}f_i$（沿 $\mathbf{f}$ 的 Lie 导数，几何来源见 \cref{ch:lie}）。把 $V=x_1^2+x_2^2$ 写成 $2x_1+2x_2$ 是\emph{梯度}不是 $\dot V$。
\quad\textbf{(ii) $\dot V=0$ 意味系统不动？}\;否。$\dot V=0$ 只表示轨迹在 $V$ 的等值面上滑行，系统可在 level set 上高速运动而 $V$ 不变（如卫星沿等势面绕行）。这一点正是下节 LaSalle 的着力处。
\end{pitfall}

% ════════════════════════════════════════════════════════════════
\section{LaSalle 不变集原理：完整证明}
\label{sec:lyap-lasalle}

\paragraph{动机：工程里 $\dot V<0$ 几乎不可能。}
机械系统的总能量 $V=T+U$ 沿轨迹的变化率是 $\dot V=-\text{耗散功率}\le0$，而耗散\emph{只在运动时}发生：当 $\dot{\mathbf{q}}=\mathbf{0}$ 时阻尼为零，于是 $\dot V=0$ 不仅在原点、在整个 $\{\dot{\mathbf{q}}=\mathbf{0}\}$ 上都成立。$\dot V$ 只是半负定，直接法仅给“稳定”。但物理直觉确信有阻尼的摆\emph{必然}停下——LaSalle 正是把这一信念严格化的工具，它是 PD 控制（\cref{sec:lyap-robot-pd}）稳定性证明的关键。

\begin{theorem}[LaSalle 不变集原理，完整版；深化 \cref{thm:ctrl-lasalle}]
\label{thm:lyap-lasalle-full}
设 $\Omega\subset\mathcal{D}$ 为紧正不变集，$V\in C^1(\Omega)$ 且 $\dot V\le0$ 于 $\Omega$。记 $E=\{\mathbf{x}\in\Omega:\dot V(\mathbf{x})=0\}$，$M$ 为 $E$ 中\emph{最大不变集}。则由 $\Omega$ 出发的每条轨迹满足 $\mathbf{x}(t)\to M\ (t\to\infty)$。特别地，若 $M=\{\mathbf{0}\}$ 则原点渐近稳定；若再 $V$ 径向无界、$\Omega=\mathbb{R}^n$ 则 GAS。
\end{theorem}

\begin{proof}
\emph{(1)} $\dot V\le0\Rightarrow V(\mathbf{x}(t))$ 单调不增，且 $V$ 在紧集 $\Omega$ 上有下界，故 $c_0=\lim_{t\to\infty}V(\mathbf{x}(t))$ 存在。
\emph{(2)} 轨迹困于紧集 $\Omega$，由 Bolzano–Weierstrass，$\omega$-极限集 $\omega(\mathbf{x}_0)=\{\mathbf{y}:\exists t_n\to\infty,\ \mathbf{x}(t_n)\to\mathbf{y}\}$ 非空、紧，且（$\omega$-极限集标准性质）\emph{不变}。
\emph{(3)} 对 $\mathbf{y}\in\omega(\mathbf{x}_0)$ 取 $t_n\to\infty,\ \mathbf{x}(t_n)\to\mathbf{y}$，由 $V$ 连续 $V(\mathbf{y})=\lim_n V(\mathbf{x}(t_n))=c_0$，故 $V\equiv c_0$ 于 $\omega(\mathbf{x}_0)$。
\emph{(4)} $\omega(\mathbf{x}_0)$ 不变且 $V$ 在其上为常数，故沿其中轨迹 $\dot V=0$，即 $\omega(\mathbf{x}_0)\subseteq E$。
\emph{(5)} $\omega(\mathbf{x}_0)$ 不变且 $\subseteq E$，故 $\subseteq M$（$E$ 中最大不变集）。
\emph{(6)} 轨迹趋于其 $\omega$-极限集：$\operatorname{dist}(\mathbf{x}(t),\omega(\mathbf{x}_0))\to0$，结合 $\omega(\mathbf{x}_0)\subseteq M$ 得 $\operatorname{dist}(\mathbf{x}(t),M)\to0$。
\end{proof}

\begin{insight}
LaSalle 的威力\emph{不}在于把 $\dot V$ 的条件放松了，而在于“不变性”是一条极强的额外约束。$E=\{\dot V=0\}$ 可能很大，但其中“系统能\emph{永远}停留”的部分（最大不变集 $M$）通常很小——往往只有原点。这像消去法：你知道嫌犯在城里（$\Omega$）、活动范围在缩小（$V$ 递减），最终他只能待在“不动也能生存”的地方（$M$）；若那地方唯一，就找到了他。
\end{insight}

\paragraph{经典应用：阻尼单摆。}
$\ddot\theta+b\dot\theta+\sin\theta=0\ (b>0)$，状态 $\dot x_1=x_2,\ \dot x_2=-\sin x_1-bx_2$。取 $V=\tfrac12 x_2^2+(1-\cos x_1)$（动能+势能），则
\begin{equation}
  \dot V=x_2(-\sin x_1-bx_2)+x_2\sin x_1=-bx_2^2\le0,
  \label{eq:lyap-pendulum-vdot}
\end{equation}
仅半负定（整条 $\{x_2=0\}$ 上为零）。$E=\{x_2=0\}$；在 $E$ 中找不变集：若 $x_2(t)\equiv0$ 则 $\dot x_2=-\sin x_1\equiv0$，即 $x_1=n\pi$。在紧水平集 $\{V<2\}$ 内（排除 $x_1=\pi$ 的鞍点），$M=\{(0,0)\}$，故原点渐近稳定。这正是 \cref{fig:lyap-lasalle-pendulum} 的图像。

\begin{figure}[H]
\centering
\begin{tikzpicture}[scale=1.0,>=Stealth,
    lvl/.style={thick, main!70!black},
    eaxis/.style={very thick, red!70!black},
    traj/.style={->, third!70!black, thick}]
  % phase axes
  \draw[->] (-3.3,0) -- (3.5,0) node[right] {$x_1=\theta$};
  \draw[->] (0,-2.3) -- (0,2.5) node[above] {$x_2=\dot\theta$};
  % E = {x2 = 0} highlighted (the x1 axis)
  \draw[eaxis] (-3.1,0) -- (3.3,0);
  \node[red!70!black, anchor=south west, font=\small] at (1.9,0.05) {$E=\{\dot V=0\}=\{x_2=0\}$};
  % energy level sets (closed curves around origin) -- explicit absolute sizes
  \draw[lvl] (0,0) ellipse (0.945 and 0.595);
  \draw[lvl] (0,0) ellipse (1.62 and 1.02);
  \draw[lvl] (0,0) ellipse (2.295 and 1.445);
  % spiraling trajectory crossing level sets inward
  \draw[traj] plot[domain=0:6.0,samples=80,variable=\t]
    ({1.9*exp(-0.13*\t)*cos(\t r)},{1.2*exp(-0.13*\t)*sin(\t r)});
  % equilibrium and saddle markers
  \fill (0,0) circle (1.6pt) node[anchor=north east, font=\small] {$M=(0,0)$};
  \fill[gray] (3.14,0) circle (1.4pt) node[anchor=south, font=\small, gray!50!black] {$(\pi,0)$ 鞍点};
  \node[main!70!black, anchor=west, font=\small] at (-3.25,1.7) {$\{V\le c\}$};
\end{tikzpicture}
\caption{阻尼单摆的 LaSalle 图像。能量水平集 $\{V\le c\}$ 为闭曲线，$\dot V=-bx_2^2$ 在整条 $x_2$ 轴（$E$）上为零；但 $E$ 中除原点外无完整轨迹（落在 $E$ 即被重力拉离），故最大不变集 $M=\{(0,0)\}$，轨迹螺旋收敛于原点。}
\label{fig:lyap-lasalle-pendulum}
\end{figure}

\begin{pitfall}[$E$ 不是 $M$；干摩擦使 LaSalle 失效]
$E=\{\dot V=0\}$ 通常远大于 $M$：上例 $E$ 是整条 $x_2$ 轴，$M$ 只有原点（与若干鞍点）。求 $M$ 必须把 $E$ 中的点\emph{代回系统方程}检查不变性，不能把 $E$ 当 $M$。另外 LaSalle 要求 $\mathbf{f}$ 足够正则（保证解唯一、$\omega$-极限集有定义）：若把阻尼改成干摩擦 $-b\operatorname{sgn}(x_2)$，$\mathbf{f}$ 在 $x_2=0$ 不连续，经典 LaSalle 不再直接适用，须用非光滑（Filippov 解）推广。
\end{pitfall}

% ════════════════════════════════════════════════════════════════
\section{指数稳定、Lyapunov 方程与间接法}
\label{sec:lyap-exp}

\paragraph{承上：从“收敛”到“多快收敛”。}
\cref{ch:lyapunov-basics} 已会算 Lyapunov 方程，\cref{ch:control} 已陈述其与 Hurwitz 的等价（\cref{thm:ctrl-lyap-eq}）。本节补两件事：一是把“二次夹”精确翻译为指数衰减率 $k,\lambda$；二是把 Lyapunov 方程放进 Sylvester$\to$Lyapunov$\to$Riccati 的层级里，揭示它正是 LQR 闭环值函数的来源（接 \cref{ch:lqr-lqg-riccati}）。在实时控制中，“最终收敛”可能是 $0.1$ 秒也可能是 $10$ 小时——指数稳定给出明确时间常数 $1/\lambda$ 来回答“扰动后多久恢复到指定精度”。

\begin{theorem}[指数稳定的二次夹刻画]
\label{thm:lyap-exp-quad}
原点指数稳定当且仅当存在 $V$ 与 $c_1,c_2,c_3>0$ 使
\begin{equation}
  c_1\|\mathbf{x}\|^2\le V(\mathbf{x})\le c_2\|\mathbf{x}\|^2,\qquad \dot V(\mathbf{x})\le-c_3\|\mathbf{x}\|^2.
  \label{eq:lyap-exp-sandwich}
\end{equation}
此时 $\|\mathbf{x}(t)\|\le\sqrt{c_2/c_1}\,\|\mathbf{x}(0)\|\,e^{-(c_3/2c_2)t}$，即 $k=\sqrt{c_2/c_1},\ \lambda=c_3/(2c_2)$。
\end{theorem}

\begin{proof}
由上夹 $V\le c_2\|\mathbf{x}\|^2$ 与 $\dot V\le-c_3\|\mathbf{x}\|^2$ 得 $\dot V\le-(c_3/c_2)V$，比较引理（见下 \cref{lem:lyap-comparison}）给 $V(\mathbf{x}(t))\le V(\mathbf{x}(0))e^{-(c_3/c_2)t}$。再用下夹 $c_1\|\mathbf{x}\|^2\le V$：$\|\mathbf{x}(t)\|\le\sqrt{V/c_1}\le\sqrt{c_2/c_1}\|\mathbf{x}(0)\|e^{-(c_3/2c_2)t}$（开方使指数减半）。
\end{proof}

\paragraph{Lyapunov 方程的积分解与对偶辨析。}
对 $\dot{\mathbf{x}}=\mathbf{A}\mathbf{x}$，取 $V=\mathbf{x}^\top\mathbf{P}\mathbf{x}$（$\mathbf{P}\succ0$）得 $\dot V=\mathbf{x}^\top(\mathbf{A}^\top\mathbf{P}+\mathbf{P}\mathbf{A})\mathbf{x}$；令其等于 $-\mathbf{x}^\top\mathbf{Q}\mathbf{x}$ 即 Lyapunov 方程 $\mathbf{A}^\top\mathbf{P}+\mathbf{P}\mathbf{A}=-\mathbf{Q}$（\cref{eq:ctrl-lyap-eq}）。$\mathbf{A}$ Hurwitz 时其唯一正定解有\emph{积分形式}
\begin{equation}
  \mathbf{P}=\int_0^\infty e^{\mathbf{A}^\top t}\,\mathbf{Q}\,e^{\mathbf{A}t}\,\mathrm dt,
  \label{eq:lyap-int-solution}
\end{equation}
收敛性与验证已在 \cref{thm:ctrl-lyap-eq} 证明给出，此处不复述。

\begin{pitfall}[把 Lyapunov 方程写成对偶式]
$\mathbf{A}^\top\mathbf{P}+\mathbf{P}\mathbf{A}=-\mathbf{Q}$（稳定性/能观 Gramian 型，$\mathbf{P}=\int e^{\mathbf{A}^\top t}\mathbf{Q}e^{\mathbf{A}t}\mathrm dt$）与\emph{对偶}方程 $\mathbf{A}\mathbf{P}+\mathbf{P}\mathbf{A}^\top=-\mathbf{Q}$（能控 Gramian 型，$\mathbf{P}=\int e^{\mathbf{A}t}\mathbf{Q}e^{\mathbf{A}^\top t}\mathrm dt$）容易写反。判稳要用前者；调用 \texttt{scipy.\allowbreak linalg.\allowbreak solve\_continuous\_lyapunov(A.\allowbreak T, -Q)} 时务必传转置。两者的对偶正是能控/能观对偶（\cref{ch:controllability-obsv}）的矩阵方程版本。
\end{pitfall}

\begin{insight}[Lyapunov 是 Riccati 的退化、是 LQR 值函数]
代数 Riccati 方程 $\mathbf{A}^\top\mathbf{P}+\mathbf{P}\mathbf{A}-\mathbf{P}\mathbf{B}\mathbf{R}^{-1}\mathbf{B}^\top\mathbf{P}+\mathbf{Q}=0$ 比 Lyapunov 方程多一个二次项 $-\mathbf{P}\mathbf{B}\mathbf{R}^{-1}\mathbf{B}^\top\mathbf{P}$——那正是“控制优化的贡献”。当 $\mathbf{B}=\mathbf{0}$（无控）Riccati 退化为 Lyapunov 方程，这给出 Sylvester$\to$Lyapunov$\to$Riccati 的层级（\cref{ch:lqr-lqg-riccati} 系统展开）。反过来，LQR 值函数 $V^\star=\mathbf{x}^\top\mathbf{P}^\star\mathbf{x}$（$\mathbf{P}^\star$ 为 CARE 解）本身就是闭环 $\dot{\mathbf{x}}=(\mathbf{A}-\mathbf{B}\mathbf{K}^\star)\mathbf{x}$ 的 Lyapunov 函数，把 Riccati 改写为闭环 Lyapunov 形式恰得 $\dot V^\star=-\mathbf{x}^\top(\mathbf{Q}+\mathbf{K}^{\star\top}\mathbf{R}\mathbf{K}^\star)\mathbf{x}<0$（\cref{thm:ctrl-lqr-lyap}）。这一联系在 CLF-QP（\cref{ch:clf-cbf}）与 MPC 终端代价（\cref{ch:mpc-advanced}）反复出现。
\end{insight}

\paragraph{间接法的三种情形与临界失效。}
设 $\mathbf{A}=\nabla\mathbf{f}(\mathbf{0})$。\cref{thm:ctrl-indirect} 已给三分法：全负实部 $\Rightarrow$ 局部指数稳定；有正实部 $\Rightarrow$ 不稳定；\emph{有零实部} $\Rightarrow$ 不可判，须用直接法。第三种失效的根源是：线性化只捕捉一阶行为，虚轴特征值处一阶近似中性稳定，真实行为由高阶项决定。典型失效例 $\dot x_1=x_2,\ \dot x_2=-x_1^3$，$\mathbf{A}=\big[\begin{smallmatrix}0&1\\0&0\end{smallmatrix}\big]$ 特征值全零——线性化完全无能；但取 $V=x_1^4/4+x_2^2/2$ 得 $\dot V=x_1^3x_2+x_2(-x_1^3)=0$，知其为保守系统、原点 Lyapunov 稳定（非渐近）。这再次说明：\cref{ch:control} 的“线性化稳定 $\Rightarrow$ 非线性全局稳定”是\emph{错}的——$\dot x=-x+x^3$ 在原点局部稳定（$\mathbf{A}=-1$）但 $|x|>1$ 不稳定。

\begin{lemma}[比较引理与 Gronwall–Bellman；Khalil Lemma 4.4]
\label{lem:lyap-comparison}
若标量 $\dot v=-\alpha(v),\ v(0)=v_0,\ \alpha\in\mathcal{K}$，则存在唯一 $\sigma\in\mathcal{KL}$ 使 $v(t)=\sigma(v_0,t)$；从而 $\dot V\le-\alpha(V)\Rightarrow V(\mathbf{x}(t))\le\sigma(V(\mathbf{x}(0)),t)$。指数情形 $\dot V\le-\lambda V\Rightarrow V(t)\le V(0)e^{-\lambda t}$。又 Gronwall–Bellman：$u(t)\le c+\int_0^t k(s)u(s)\,\mathrm ds\Rightarrow u(t)\le c\exp\!\big(\int_0^t k(s)\,\mathrm ds\big)$。
\end{lemma}

\noindent 比较引理把 Lyapunov 不等式\emph{量化}为显式衰减包络，是反步法、自适应控制、观测器收敛率分析的标准“流水线”末端工具（\cref{sec:lyap-iss,sec:lyap-backstepping} 反复用到）。工程范式四步：构造 $V$ $\to$ 算 $\dot V$ $\to$ 用 Young/Cauchy–Schwarz 压成 $-\alpha(V)+\sigma(\|\mathbf{u}\|)$ $\to$ 应用比较引理得收敛率。

% ════════════════════════════════════════════════════════════════
\section{Chetaev 不稳定性定理与 Van der Pol 振子}
\label{sec:lyap-chetaev}

\paragraph{动机：怎么\emph{证明}不稳定。}
直接法告诉你“找到 $V$ 使 $\dot V\le0$ 则稳定”。但找不到这样的 $V$，是系统不稳定还是你构造能力不足？逆定理（\cref{sec:lyap-converse}）将说明“找不到”不能证不稳定。Chetaev 定理给出\emph{正面}判定不稳定的工具——找一个“反 Lyapunov 函数”，在原点旁的一个锥形区域内 $V$ 正且 $\dot V$ 正。

\begin{theorem}[Chetaev 不稳定性定理]
\label{thm:lyap-chetaev}
若存在 $C^1$ 函数 $V$ 与原点任意邻域中的开集 $U$，使 \textbf{(1)} $V>0$ 于 $U$；\textbf{(2)} $V=0$ 于 $\partial U\cap\mathcal{D}$；\textbf{(3)} $\dot V>0$ 于 $U$，则原点\emph{不稳定}。
\end{theorem}

\noindent 几何直觉：在原点旁的锥形区域内 $V$ 正且沿轨迹递增，轨迹一旦进入就像被推离的水球越走越远，无法保持在原点附近。应用：倒立摆上平衡点 $(\theta=\pi,\dot\theta=0)$ 的不稳定、机械臂奇异位形分析。最简验证 $\dot x_1=x_2,\dot x_2=x_1$：取 $V=x_1x_2,\ U=\{x_1x_2>0\}$，则 $\dot V=x_2^2+x_1^2>0$，故原点不稳定（与“$\mathbf{A}$ 有正实部特征值 $+1$”一致）。

\paragraph{Van der Pol 振子：Lyapunov 给出的全局定性图像。}
$\dot x_1=x_2,\ \dot x_2=-x_1+\mu(1-x_1^2)x_2\ (\mu>0)$。取 $V=\tfrac12(x_1^2+x_2^2)$ 得
\begin{equation}
  \dot V=\mu(1-x_1^2)x_2^2.
  \label{eq:lyap-vdp-vdot}
\end{equation}
在带状区 $|x_1|<1$（且 $x_2\neq0$）$\dot V>0$（能量增），取 Chetaev 集 $U=\{x_1^2<1,\ x_2\neq0\}$ 即证原点不稳定；在 $|x_1|>1$ 的远处 $\dot V<0$（能量减），轨迹被吸引回来。\textbf{原点排斥 + 远处吸引 $\Rightarrow$ 存在极限环}——这是非线性振荡的经典图像（\cref{fig:lyap-vdp}）。

\begin{figure}[H]
\centering
\begin{tikzpicture}[scale=1.0,>=Stealth]
  \draw[->] (-3.2,0) -- (3.4,0) node[right] {$x_1$};
  \draw[->] (0,-2.6) -- (0,2.7) node[above] {$x_2$};
  % the strip |x1|<1 where Vdot>0
  \begin{scope}
    \fill[red!8] (-1,-2.5) rectangle (1,2.6);
  \end{scope}
  \draw[red!50,dashed] (-1,-2.5) -- (-1,2.6);
  \draw[red!50,dashed] (1,-2.5) -- (1,2.6);
  \node[red!60!black, anchor=south, font=\small] at (0,2.0) {$\dot V>0$};
  \node[main!60!black, anchor=south west, font=\small] at (2.0,1.9) {$\dot V<0$};
  % limit cycle (closed orbit)
  \draw[very thick, third!70!black] (0,0) ellipse (2.1 and 1.9);
  \node[third!70!black, anchor=west, font=\small] at (2.12,0.4) {极限环};
  % outer trajectory spiraling inward to the cycle
  \draw[->, gray!60!black] plot[domain=0:4.2,samples=70,variable=\t]
    ({(2.9-0.18*\t)*cos(60*\t)},{(2.6-0.16*\t)*sin(60*\t)});
  % inner trajectory spiraling outward
  \draw[->, gray!60!black] plot[domain=0:3.6,samples=70,variable=\t]
    ({(0.25+0.16*\t)*cos(60*\t)},{(0.22+0.15*\t)*sin(60*\t)});
  \fill (0,0) circle (1.5pt) node[anchor=north east, font=\small] {不稳定原点};
\end{tikzpicture}
\caption{Van der Pol 振子（$\mu>0$）的相图。带状区 $|x_1|<1$ 内 $\dot V>0$（原点被排斥，Chetaev），远处 $\dot V<0$（轨迹被吸引）；内外轨迹都卷向同一条稳定极限环。}
\label{fig:lyap-vdp}
\end{figure}

\begin{pitfall}[Poincaré–Bendixson 给\emph{存在性}，不给唯一性]
极限环的\emph{存在性}由 Poincaré–Bendixson 定理保证（平面系统、紧不变环域内无平衡点则有周期轨）。但 PB 定理\emph{不}保证唯一性——Van der Pol 极限环的唯一性须另由 Liénard 定理 / Levinson–Smith 定理给出。把“PB $\Rightarrow$ 唯一极限环”写进证明是常见的越界。还须注意：经典 Lyapunov 以\emph{平衡点}为中心，极限环是周期轨道而非平衡点，分析其稳定性要升级工具（Poincaré 映射 + 特征乘子，或 \cref{sec:lyap-outlook} 的 Contraction）。
\end{pitfall}

% ════════════════════════════════════════════════════════════════
\section{Lyapunov 函数的构造方法}
\label{sec:lyap-construct-d}

\paragraph{Lyapunov 方法的“阿喀琉斯之踵”。}
直接法的逻辑优美，却有一个根本难题：\emph{$V$ 从哪来？}定理只说“若存在则稳定”，不说怎么找。\cref{ch:lyapunov-basics} 已系统讲过\emph{能量法、二次型、Krasovskii 法、变量梯度法}的计算细节，本节不重复其推导，只\textbf{补两件}：(i) Krasovskii 法“全局”结论一个常被略去的条件；(ii) 两种把构造\emph{计算化}的现代方法（SOS、Neural Lyapunov）的概念与定理。

\paragraph{能量法与二次型（回顾接口）。}
机械系统取 $V=\tfrac12\dot{\mathbf{q}}^\top\mathbf{M}(\mathbf{q})\dot{\mathbf{q}}+U(\mathbf{q})$，借 $\dot{\mathbf{M}}-2\mathbf{C}$ 反对称（来自拉格朗日结构，几何来源见 \cref{ch:rigid_body}，空间向量代数将在 \cref{ch:spatial-dynamics} 正式证）得 $\dot V=-\dot{\mathbf{q}}^\top\mathbf{D}\dot{\mathbf{q}}\le0$，配 LaSalle 判渐近稳定（\cref{sec:lyap-robot-pd} 详证）。局部分析则取二次型 $V=\mathbf{x}^\top\mathbf{P}\mathbf{x}$ 解 Lyapunov 方程，$\dot V=-\mathbf{x}^\top\mathbf{Q}\mathbf{x}+O(\|\mathbf{x}\|^3)$ 给局部指数稳定。两法均在 \cref{ch:lyapunov-basics} 算透。

\begin{pitfall}[Krasovskii 法的“全局”要补径向无界]
Krasovskii 法（\cref{ch:lyapunov-basics}）取 $V=\|\mathbf{f}(\mathbf{x})\|^2$，由雅可比 $\mathbf{J}(\mathbf{x})=\nabla\mathbf{f}$ 满足 $\mathbf{J}+\mathbf{J}^\top\preceq-\alpha\mathbf{I}\ (\alpha>0)$ 得 $\dot V\le-\alpha\|\mathbf{f}\|^2<0\ (\mathbf{x}\neq\mathbf{0})$，从而\emph{渐近稳定}。但要升级到\textbf{全局}（GAS），还须 $V=\|\mathbf{f}(\mathbf{x})\|^2$ 本身\emph{径向无界}，即 $\|\mathbf{f}(\mathbf{x})\|\to\infty$ 当 $\|\mathbf{x}\|\to\infty$——这是经典 Krasovskii–Barbashin 全局定理的标准附加条件。仅“一致负定 + $\mathcal{D}=\mathbb{R}^n$”\emph{不}自动蕴含 $\|\mathbf{f}\|$ 径向无界（否则 \cref{thm:lyap-gas} 的紧性会失守）。写法：$\mathbf{J}+\mathbf{J}^\top\preceq-\alpha\mathbf{I}\Rightarrow$ 局部 AS；GAS 须另加“且 $\|\mathbf{f}(\mathbf{x})\|\to\infty$”一句。
\end{pitfall}

\paragraph{方法五：SOS 多项式搜索（概念）。}
对多项式系统，把 $V$ 参数化、把 Lyapunov 条件转化为半定规划（SDP）。多项式 $p$ 称 \emph{Sum of Squares (SOS)}，若 $p=\sum_i q_i^2$（$q_i$ 为多项式）；SOS $\Rightarrow p\ge0$（充分，一般有 gap）。算法：参数化 $V(\mathbf{x})=\mathbf{m}(\mathbf{x})^\top\mathbf{P}\mathbf{m}(\mathbf{x})$（$\mathbf{m}$ 为单项式基，$\mathbf{P}\succeq0$），要求 $V-\epsilon\|\mathbf{x}\|^2$ 与 $-\dot V-\epsilon\|\mathbf{x}\|^2$ 均为 SOS，用 SDP 求解器解 $\mathbf{P}$（SDP 背景见 \cref{ch:nlopt}）。优势是完全自动、可同时优化吸引域；局限是限于多项式系统、随维数/阶数计算量快升、SOS 的 gap。Parrilo 2000 博士论文\cite{parrilo2000structured} 奠定了这一凸化框架（工具链 API 不在本书范围）。

\paragraph{方法六：Neural Lyapunov（概念）。}
用神经网络 $V_\theta(\mathbf{x})$ 参数化 Lyapunov 函数，训练目标为 Lyapunov 条件的违反量 $\max(0,-V_\theta+\epsilon\|\mathbf{x}\|^2)+\max(0,\dot V_\theta+\alpha V_\theta)$；再用 SMT 求解器（dReal）做\emph{全域}反例搜索（falsifier），迭代到找不到反例为止。Chang–Roohi–Gao 2019\cite{chang2019neural} 的 learner+falsifier 框架借 dReal 的 $\delta$-complete 判定给出形式化保证。这与 SOS 同属“把构造变成搜索”的范式，其可行性由逆定理（\cref{sec:lyap-converse}）保证——目标 $V$ 必然存在。

\begin{table}[H]\centering\small
\caption{Lyapunov 函数构造方法对比}
\label{tab:lyap-construct}
\begin{tabular}{lllll}
\toprule
方法 & 适用范围 & 自动化 & 全局性 & 计算量 \\
\midrule
能量法 & 机械系统 & 手工 & 局部$\to$全局需 LaSalle & 极低 \\
二次型+Lyap 方程 & 线性/局部 & 半自动 & 局部 & 低 \\
Krasovskii & $\mathbf{J}+\mathbf{J}^\top\prec0$ & 检验型 & 全局（须径向无界） & 低 \\
变量梯度 & 低维 & 手工 & 视情况 & 中 \\
SOS & 多项式系统 & 全自动 & 可优化 ROA & 高 \\
Neural & 通用 & 全自动 & 需验证 & 很高 \\
\bottomrule
\end{tabular}
\end{table}

\begin{pitfall}[“找不到 $V$” 与 “SOS 不可行” 都不证不稳定]
逆定理（下节）保证稳定系统一定有 $V$；找不到只意味构造技巧或算力不足，\emph{不能}据此判不稳定。同理 SOS 在低阶 $V$ 上声称“不可行”，只是搜索空间太小——应逐步增阶，而非断言系统不稳定。
\end{pitfall}

% ════════════════════════════════════════════════════════════════
\section{逆定理：稳定性 \texorpdfstring{$\Leftrightarrow$}{<=>} Lyapunov 函数存在}
\label{sec:lyap-converse}

\paragraph{动机：Lyapunov 方法“完备”吗？}
每次构造失败你都会问：是否存在某些稳定系统\emph{根本没有} Lyapunov 函数？若是，方法有内在局限；若否，它原则上\emph{万能}。逆定理给出后者：任何稳定系统都存在 $V$，找不到只是人（或机器）的构造能力不足。

\begin{theorem}[逆定理三阶梯：Massera / Kurzweil / Lin–Sontag–Wang]
\label{thm:lyap-converse}
\textbf{(Massera 1956\cite{massera1956contributions})}\;若 $\dot{\mathbf{x}}=\mathbf{f}(\mathbf{x})$ 原点渐近稳定，则存在 $C^\infty$ 函数 $V$ 使 $\dot V<0\ (\mathbf{x}\neq\mathbf{0})$。
\textbf{(Kurzweil 1956\cite{kurzweil1956inversion})}\;若原点\emph{全局}渐近稳定，则存在\emph{全局定义}的 $C^\infty$ Lyapunov 函数。
\textbf{(Lin–Sontag–Wang 1996\cite{lin1996smooth})}\;即便对带时变扰动的非自治系统，只要在 $\mathcal{KL}$ 意义下（鲁棒地）稳定，就存在\emph{光滑} ISS-Lyapunov 函数。
\end{theorem}

\noindent Massera 的构造思路（仅概述）：由 $\mathcal{KL}$ 衰减 $\|\mathbf{x}(t)\|\le\beta(\|\mathbf{x}_0\|,t)$ 沿轨迹积分 $V(\mathbf{x})=\int_0^\infty g(\|\boldsymbol{\phi}(t;\mathbf{x})\|)\,\mathrm dt$，其中 $\boldsymbol{\phi}(t;\mathbf{x})$ 是从 $\mathbf{x}$ 出发的流、$g$ 取使积分收敛且 $V$ 光滑的核。

\begin{insight}[逆定理的哲学价值]
逆定理把 Lyapunov 方法从“有时成功的充分条件”提升为“原则上总能成功的\emph{充要}判据”：证明稳定性\emph{不存在}比 Lyapunov 更强的方法。对 SOS / Neural Lyapunov，它保证了\emph{搜索目标的存在性}——你不是在搜一个可能不存在的东西。倘若逆定理不成立（存在稳定却无 $V$ 的系统），则 SOS 永不可能找到证书、Neural 训练永不收敛，你将被迫去解 ODE——对复杂系统不可能。所幸这种情形不会发生。
\end{insight}

\begin{pitfall}[存在性 $\neq$ 构造性]
逆定理是\emph{存在性}定理。Massera/Kurzweil 给的积分式 $V=\int_0^\infty g(\|\boldsymbol{\phi}(t;\mathbf{x})\|)\mathrm dt$ 要算它就得知道精确流 $\boldsymbol{\phi}$（即解方程），又回到起点。其价值是哲学性的：保证搜索方向正确，而非交付一个可计算的 $V$。另外非光滑系统也有逆定理（Clarke–Ledyaev–Stern / Teel–Praly 的 Filippov 解推广），故“非光滑系统没有逆定理”亦是误区。
\end{pitfall}

% ════════════════════════════════════════════════════════════════
\section{ISS：输入到状态稳定}
\label{sec:lyap-iss}

\paragraph{动机：真实系统从不封闭。}
至此都是自治系统 $\dot{\mathbf{x}}=\mathbf{f}(\mathbf{x})$。但真实机器人总有外部扰动（电机噪声、量化、观测器误差、地面摩擦扰动、风/碰撞），系统变为 $\dot{\mathbf{x}}=\mathbf{f}(\mathbf{x},\mathbf{u})$。此时经典 Lyapunov 不再适用——持续扰动会阻止轨迹趋零，你不能要求 $\mathbf{x}\to\mathbf{0}$。Sontag 1989\cite{sontag1989smooth} 提出 ISS（Input-to-State Stability）正是为定量回答：\emph{有持续扰动时，轨迹偏离原点多远？偏移量与扰动大小是什么定量关系？}这是把 \cref{sec:lyap-comparison} 的比较函数语言用到鲁棒性上的第一个大成果，也是 CLF/CBF（\cref{ch:clf-cbf}）的直接前置。

\begin{definition}[ISS；Sontag 1989]
\label{def:lyap-iss}
$\dot{\mathbf{x}}=\mathbf{f}(\mathbf{x},\mathbf{u})$ 称 ISS，若存在 $\beta\in\mathcal{KL},\ \gamma\in\mathcal{K}$ 使对所有可测有界 $\mathbf{u}$ 与所有初值 $\mathbf{x}_0$：
\begin{equation}
  \|\mathbf{x}(t)\|\le\beta(\|\mathbf{x}_0\|,t)+\gamma\Big(\sup_{0\le\tau\le t}\|\mathbf{u}(\tau)\|\Big).
  \label{eq:lyap-iss-def}
\end{equation}
\end{definition}

\noindent 物理含义：$\beta(\|\mathbf{x}_0\|,t)$ 是初值影响随时间衰减（如无扰渐近稳定），$\gamma(\|\mathbf{u}\|_\infty)$ 把扰动影响“限住”——扰动越大偏离越大，但有界。类比弹簧-阻尼受外力：恢复力保证不无限偏离、阻尼保证暂态衰减，$\gamma$ 是“柔度”，越小系统越“硬”、越鲁棒。注意 ISS 不蕴含“扰动减小则趋零”：若 $\mathbf{u}$ 只是变小而非消失，状态趋向 $\gamma(\limsup\|\mathbf{u}\|)$ 邻域而非原点；唯 $\mathbf{u}\equiv\mathbf{0}$ 时退化为 $\|\mathbf{x}(t)\|\le\beta(\|\mathbf{x}_0\|,t)\to0$。

\begin{theorem}[ISS 等价刻画；Sontag–Wang]
\label{thm:lyap-iss-char}
对正则右端 $\mathbf{f}$，以下等价：\textbf{(1)} 系统 ISS；\textbf{(2)} 0-GAS（零输入下 GAS）\emph{且}具\emph{渐近增益}（AG：$\exists\gamma\in\mathcal{K}$ 使 $\limsup_{t\to\infty}\|\mathbf{x}(t)\|\le\gamma(\|\mathbf{u}\|_\infty)$）；\textbf{(3)} 存在（光滑）ISS-Lyapunov 函数\cite{sontagwang1995characterizations}。\footnote{此处 (2) 的判据须取\emph{渐近增益} AG，而非仅“有界输入有界状态”（BIBS/UBIBS）：ISS 蕴含 BIBS，但反向不真——0-GAS$+$BIBS 一般\emph{不}足以推出 ISS，缺的正是 AG 这条“偏移量由扰动渐近幅值锁定”的性质。AG 刻画见 Sontag \& Wang, \emph{New Characterizations of Input-to-State Stability}, IEEE TAC 41(9):1283–1294, 1996；(3) 的 Lyapunov 等价见 \cite{sontagwang1995characterizations}（Syst.\ Control Lett.\ 24, 1995）。}
\end{theorem}

\begin{definition}[ISS-Lyapunov 函数]
\label{def:lyap-iss-lyap}
$V:\mathbb{R}^n\to\mathbb{R}$ 称 ISS-Lyapunov 函数，若 $\exists\alpha_1,\alpha_2\in\mathcal{K}_\infty,\ \alpha_3,\chi\in\mathcal{K}$ 使 $\alpha_1(\|\mathbf{x}\|)\le V\le\alpha_2(\|\mathbf{x}\|)$ 且
\begin{equation}
  \|\mathbf{x}\|\ge\chi(\|\mathbf{u}\|)\ \Rightarrow\ \dot V\le-\alpha_3(\|\mathbf{x}\|)\quad(\text{蕴含型}),
  \label{eq:lyap-iss-implication}
\end{equation}
等价的\emph{耗散型}：$\dot V(\mathbf{x},\mathbf{u})\le-\alpha(\|\mathbf{x}\|)+\sigma(\|\mathbf{u}\|),\ \alpha,\sigma\in\mathcal{K}$。
\end{definition}

\noindent 直觉：ISS-Lyapunov 说“只要状态够大（$\|\mathbf{x}\|\ge\chi(\|\mathbf{u}\|)$）能量就下降”，于是系统最终被拉入由扰动大小决定的球 $\|\mathbf{x}\|\le\chi(\|\mathbf{u}\|_\infty)$ 内。构造四步套路：选 $u=0$ 的 $V$（如 $\mathbf{x}^\top\mathbf{P}\mathbf{x}$）$\to$ 沿 $\mathbf{f}(\mathbf{x},\mathbf{u})$ 算 $\dot V$ $\to$ 用 Young 不等式 $ab\le\tfrac{a^2}{2\epsilon}+\tfrac{\epsilon b^2}{2}$ 拆交叉项 $\to$ 验耗散型条件。

\begin{example}[标量 ISS 范例]
\label{ex:lyap-iss}
$\dot x=-x^3+xu$，取 $V=\tfrac12 x^2$，则 $\dot V=-x^4+x^2u$。用 Young 不等式 $x^2u\le x^2|u|\le\tfrac12 x^4+\tfrac12 u^2$ 得
\begin{equation}
  \dot V\le-x^4+\tfrac12 x^4+\tfrac12 u^2=-\tfrac12 x^4+\tfrac12 u^2,
  \label{eq:lyap-iss-example}
\end{equation}
即耗散型 $\alpha(r)=\tfrac12 r^4,\ \sigma(r)=\tfrac12 r^2$。由 $\alpha(\chi(r))=\sigma(r)$ 解出 ISS 增益 $\chi(r)=r^{1/2}$。
\end{example}

\paragraph{ISS 小增益定理。}
现代机器人控制充满\emph{互联}：控制器以估计状态为输入、观测器以控制输出为输入；分布式多机经网络耦合。小增益定理给出互联稳定性的判据。

\begin{theorem}[ISS 小增益定理；Jiang–Teel–Praly 1994]
\label{thm:lyap-small-gain}
两子系统反馈互联 $\dot{\mathbf{x}}_1=\mathbf{f}_1(\mathbf{x}_1,\mathbf{x}_2),\ \dot{\mathbf{x}}_2=\mathbf{f}_2(\mathbf{x}_1,\mathbf{x}_2)$。若 $\Sigma_1$ 以 $\mathbf{x}_2$ 为输入 ISS（增益 $\gamma_1$）、$\Sigma_2$ 以 $\mathbf{x}_1$ 为输入 ISS（增益 $\gamma_2$），且\emph{小增益条件}
\begin{equation}
  \gamma_1\circ\gamma_2(r)<r,\qquad\forall r>0,
  \label{eq:lyap-small-gain}
\end{equation}
成立，则整体 $(\Sigma_1,\Sigma_2)$ ISS\cite{jiang1994smallgain}。
\end{theorem}

\begin{insight}[非线性小增益 $\leftrightarrow$ 频域 $H_\infty$ 小增益]
每个子系统的输出被另一个“放大”后反馈；若合成增益 $\gamma_1\circ\gamma_2$ 小于恒等（衰减），扰动在环路中越来越小，整体稳定。退化到线性时增益是 $L_\infty$ 增益 $\gamma_i(r)=\|G_i\|_\infty\cdot r$，\cref{eq:lyap-small-gain} 化为经典 $\|G_1\|_\infty\|G_2\|_\infty<1$——这正是鲁棒控制判互联稳定性的核心工具。\textbf{分工}：本章 ISS 小增益是\emph{非线性时域}的互联判据；其\emph{频域线性}版本（$\|G_1\|_\infty\|G_2\|_\infty<1$ 与回路成形）属 \cref{ch:robust-hinf}，两者互引、不重复推导线性退化。
\end{insight}

\noindent 与 ISS 平行的弱化概念是 \emph{iISS}（积分 ISS，Sontag 1998\cite{sontag1998mathematical}）：$\|\mathbf{x}(t)\|\le\beta(\|\mathbf{x}_0\|,t)+\int_0^t\gamma(\|\mathbf{u}(\tau)\|)\mathrm d\tau$。iISS 只要求“积分有限的输入 $\Rightarrow$ 有界状态”，比 ISS 弱，对双线性系统等更适用。小增益定理的力量在于：\emph{一旦知道增益就能直接判互联稳定}，省去为整体重构 Lyapunov 函数——但计算非线性增益 $\gamma_i$ 本身往往就要找 ISS-Lyapunov 函数，并不轻松。

% ════════════════════════════════════════════════════════════════
\section{级联系统稳定性}
\label{sec:lyap-cascade}

\paragraph{动机：层级控制即级联。}
现代机器人控制几乎都是层级结构：高层规划输出参考、底层控制跟踪。数学上对应\emph{级联系统} $\dot{\mathbf{x}}_1=\mathbf{f}_1(\mathbf{x}_1)$（上层，不受下层影响）、$\dot{\mathbf{x}}_2=\mathbf{f}_2(\mathbf{x}_1,\mathbf{x}_2)$（下层，受上层驱动）。核心问题：上层单独 GAS、下层零输入 $\dot{\mathbf{x}}_2=\mathbf{f}_2(\mathbf{0},\mathbf{x}_2)$ 也 GAS，整体是否 GAS？

\begin{theorem}[级联稳定性；Sontag 1989]
\label{thm:lyap-cascade}
若 \textbf{(1)} $\dot{\mathbf{x}}_1=\mathbf{f}_1(\mathbf{x}_1)$ 原点 GAS；\textbf{(2)} $\dot{\mathbf{x}}_2=\mathbf{f}_2(\mathbf{0},\mathbf{x}_2)$ 原点 GAS；\textbf{(3)} $\dot{\mathbf{x}}_2=\mathbf{f}_2(\mathbf{x}_1,\mathbf{x}_2)$ 以 $\mathbf{x}_1$ 为输入 ISS（或满足适当增长条件），则级联原点 $(\mathbf{0},\mathbf{0})$ GAS。
\end{theorem}

\noindent 直觉：上层最终趋零（条件 1），对下层的驱动 $\mathbf{x}_1(t)\to\mathbf{0}$；驱动消失后下层也趋零（条件 2）；条件 3 保证过渡期下层不“爆炸”。

\begin{pitfall}[缺条件 3 会爆炸]
仅条件 1+2 \emph{不}够。反例 $\dot x_1=-x_1,\ \dot x_2=-x_2+x_1x_2^2$：上层 GAS（$x_1$ 指数衰减）、下层零输入 GAS（$\dot x_2=-x_2$）；但若 $x_2(0)$ 很大，在 $x_1$ 衰减之前项 $x_1x_2^2$ 可令 $x_2$ 有限时间爆炸。须 ISS 类增长条件排除。
\end{pitfall}

\begin{table}[H]\centering\small
\caption{机器人中的典型级联结构}
\label{tab:lyap-cascade}
\begin{tabular}{llll}
\toprule
层级 & 上层 $\Sigma_1$ & 下层 $\Sigma_2$ & 级联条件 \\
\midrule
MPC + WBC & MPC 输出 CoM 轨迹 & WBC 跟踪关节力矩 & MPC 稳定 + WBC ISS \\
落脚规划 + 腿控 & Raibert 落脚 & 单腿摆动控制 & 落脚稳定 + 腿控 ISS \\
感知 + 控制 & SLAM 状态估计 & 基于估计的控制器 & SLAM 收敛 + 控制器对估计误差 ISS \\
\bottomrule
\end{tabular}
\end{table}

\noindent “MPC（$50$ Hz）+ PD（$1$ kHz）”天然满足级联 ISS：MPC 输出是慢变有界参考，PD 对其 ISS。这条层级脉络在 \cref{ch:control} 足式概览与 \cref{ch:force-wbc}（WBC）里反复出现；MPC 终端代价作为闭环 Lyapunov 函数（\cref{thm:ctrl-mpc-main}）则把本章 \cref{sec:lyap-direct,sec:lyap-lasalle} 直接接入 \cref{ch:mpc-advanced}。

% ════════════════════════════════════════════════════════════════
\section{反步法 Backstepping}
\label{sec:lyap-backstepping}

\paragraph{从分析到设计。}
至此 Lyapunov 主要用于\emph{分析}——给定系统判稳。它同样能用于\emph{设计}——递归构造控制律与 Lyapunov 函数。反步法是这一思想最经典的实现，适用于“可从后往前逐步镇定”的\emph{严格反馈}结构
\begin{equation}
  \dot x_1=f_1(x_1)+g_1(x_1)x_2,\qquad \dot x_2=f_2(x_1,x_2)+g_2(x_1,x_2)u.
  \label{eq:lyap-strict-feedback}
\end{equation}
关键观察：若把 $x_2$ 视作“虚拟控制”，能选 $x_2=\phi(x_1)$ 使第一式稳定，则真控制 $u$ 只需让 $x_2$ 跟踪 $\phi(x_1)$。

\begin{theorem}[反步法递归设计；Krstić–Kanellakopoulos–Kokotović 1995]
\label{thm:lyap-backstepping}
对 \cref{eq:lyap-strict-feedback}：\textbf{(1)} 取 $V_1(x_1)$ 与虚拟控制 $\phi(x_1)$（$\phi(0)=0$）使 $\dot V_1=\nabla V_1\cdot[f_1+g_1\phi]\le-W_1(x_1)<0$；\textbf{(2)} 定义误差 $z_2=x_2-\phi(x_1)$；\textbf{(3)} 增广 $V_2=V_1+\tfrac12 z_2^2$；\textbf{(4)} 算 $\dot V_2=\dot V_1+z_2(f_2+g_2u-\frac{\partial\phi}{\partial x_1}\dot x_1)$，选 $u$ 使 $\dot V_2\le-W_1(x_1)-W_2(z_2)<0$。则闭环原点渐近稳定\cite{krstic1995nonlinear}。
\end{theorem}

\begin{example}[二阶积分器的反步设计]
\label{ex:lyap-backstepping}
$\dot x_1=x_2,\ \dot x_2=u$。\textbf{(1)} 取 $V_1=x_1^2/2$，虚拟控制 $\phi=-c_1x_1\ (c_1>0)$ 给 $\dot V_1=-c_1x_1^2<0$。\textbf{(2)} $z_2=x_2+c_1x_1$。\textbf{(3)} $V_2=x_1^2/2+z_2^2/2$。\textbf{(4)} 代入展开
\begin{align}
  \dot V_2&=x_1x_2+z_2(u+c_1x_2)=x_1(z_2-c_1x_1)+z_2\big(u+c_1(z_2-c_1x_1)\big)\notag\\
  &=-c_1x_1^2+x_1z_2+z_2u+c_1z_2^2-c_1^2x_1z_2.
  \label{eq:lyap-backstep-expand}
\end{align}
选 $u=-(1-c_1^2)x_1-(c_1+c_2)z_2\ (c_2>0)$ 消去交叉项 $x_1z_2$，得 $\dot V_2=-c_1x_1^2-c_2z_2^2<0$。
\end{example}

\begin{insight}[反步法 $\leftrightarrow$ 外环-内环]
反步法的层级天然对应机器人“外环（慢）位置/姿态控制——虚拟控制是期望速度/角速度”与“内环（快）速度/角速度跟踪——用真实执行器力矩”。这与 \cref{sec:lyap-cascade} 的级联、\cref{sec:lyap-robot-pd} 的 PD 是同一层级思想的三种面孔。局限是\emph{复杂度爆炸}：每反推一层都要对 $\phi$ 求一次导（“项数膨胀”），高维系统的解析 $u$ 迅速变得不可手算——这催生了 dynamic surface control 等工程化变体。
\end{insight}

% ════════════════════════════════════════════════════════════════
\section{旗舰应用：机器人 PD+重力补偿稳定性}
\label{sec:lyap-robot-pd}

\paragraph{把全章工具拧成一个证明。}
本节是 LaSalle（\cref{sec:lyap-lasalle}）、能量法（\cref{sec:lyap-construct-d}）、$\dot{\mathbf{M}}-2\mathbf{C}$ 反对称（\cref{ch:rigid_body}）在机器人上的合流，也是 \cref{ch:control} 关节空间 PD（\cref{sec:ctrl-robot-pd}）的严格补证。

\paragraph{系统与控制律。}
$n$ 自由度机器人 $\mathbf{M}(\mathbf{q})\ddot{\mathbf{q}}+\mathbf{C}(\mathbf{q},\dot{\mathbf{q}})\dot{\mathbf{q}}+\mathbf{g}(\mathbf{q})=\boldsymbol{\tau}$，$\mathbf{M}\succ0$。\emph{PD+重力补偿}控制律
\begin{equation}
  \boldsymbol{\tau}=\mathbf{g}(\mathbf{q})-\mathbf{K}_p(\mathbf{q}-\mathbf{q}_d)-\mathbf{K}_d\dot{\mathbf{q}},\qquad\mathbf{K}_p,\mathbf{K}_d\succ0.
  \label{eq:lyap-pd-law}
\end{equation}
令 $\mathbf{e}=\mathbf{q}-\mathbf{q}_d$，闭环 $\mathbf{M}(\mathbf{q})\ddot{\mathbf{e}}+\mathbf{C}(\mathbf{q},\dot{\mathbf{q}})\dot{\mathbf{e}}+\mathbf{K}_d\dot{\mathbf{e}}+\mathbf{K}_p\mathbf{e}=\mathbf{0}$。

\begin{theorem}[PD+重力补偿的渐近稳定；Takegaki–Arimoto 1981]
\label{thm:lyap-robot-pd}
对 set-point 调节（$\dot{\mathbf{q}}_d=\mathbf{0}$），控制律 \cref{eq:lyap-pd-law} 使闭环误差 $(\mathbf{e},\dot{\mathbf{e}})=(\mathbf{0},\mathbf{0})$ 渐近稳定；当 $\mathbf{K}_p$ 充分大（主导重力梯度界）时为\emph{全局}渐近稳定\cite{takegaki1981feedback}。
\end{theorem}

\begin{proof}
取总能量型 Lyapunov 函数
\begin{equation}
  V=\tfrac12\dot{\mathbf{e}}^\top\mathbf{M}(\mathbf{q})\dot{\mathbf{e}}+\tfrac12\mathbf{e}^\top\mathbf{K}_p\mathbf{e},
  \label{eq:lyap-pd-V}
\end{equation}
由 $\mathbf{M}\succ0,\ \mathbf{K}_p\succ0$ 知 $V>0\ ((\mathbf{e},\dot{\mathbf{e}})\neq\mathbf{0})$。求导并代入闭环 $\mathbf{M}\ddot{\mathbf{e}}=-\mathbf{C}\dot{\mathbf{e}}-\mathbf{K}_d\dot{\mathbf{e}}-\mathbf{K}_p\mathbf{e}$：
\begin{align}
  \dot V&=\dot{\mathbf{e}}^\top\mathbf{M}\ddot{\mathbf{e}}+\tfrac12\dot{\mathbf{e}}^\top\dot{\mathbf{M}}\dot{\mathbf{e}}+\mathbf{e}^\top\mathbf{K}_p\dot{\mathbf{e}}
   =\dot{\mathbf{e}}^\top(-\mathbf{C}\dot{\mathbf{e}}-\mathbf{K}_d\dot{\mathbf{e}}-\mathbf{K}_p\mathbf{e})+\tfrac12\dot{\mathbf{e}}^\top\dot{\mathbf{M}}\dot{\mathbf{e}}+\mathbf{e}^\top\mathbf{K}_p\dot{\mathbf{e}}\notag\\
  &=\dot{\mathbf{e}}^\top\big(\tfrac12\dot{\mathbf{M}}-\mathbf{C}\big)\dot{\mathbf{e}}-\dot{\mathbf{e}}^\top\mathbf{K}_d\dot{\mathbf{e}}.
  \label{eq:lyap-pd-Vdot}
\end{align}
关键性质：$\dot{\mathbf{M}}-2\mathbf{C}$ 反对称（拉格朗日结构，\cref{ch:rigid_body}），故 $\dot{\mathbf{e}}^\top(\tfrac12\dot{\mathbf{M}}-\mathbf{C})\dot{\mathbf{e}}=0$，于是
\begin{equation}
  \dot V=-\dot{\mathbf{e}}^\top\mathbf{K}_d\dot{\mathbf{e}}\le0\quad(\text{半负定}).
  \label{eq:lyap-pd-Vdot-final}
\end{equation}
半负定，故须 LaSalle（\cref{thm:lyap-lasalle-full}）。$E=\{\dot V=0\}=\{\dot{\mathbf{e}}=\mathbf{0}\}$；在 $E$ 中找不变集：$\dot{\mathbf{e}}(t)\equiv\mathbf{0}\Rightarrow\ddot{\mathbf{e}}=\mathbf{0}$，代入闭环得 $\mathbf{K}_p\mathbf{e}=\mathbf{0}\Rightarrow\mathbf{e}=\mathbf{0}$。故 $M=\{(\mathbf{0},\mathbf{0})\}$，由 LaSalle $(\mathbf{e},\dot{\mathbf{e}})\to(\mathbf{0},\mathbf{0})$，即 $\mathbf{q}\to\mathbf{q}_d,\ \dot{\mathbf{q}}\to\mathbf{0}$。$\mathbf{K}_p$ 充分大时 $V$ 在全局唯一极小，得 GAS。
\end{proof}

\begin{pitfall}[为何必须重力补偿；$\dot{\mathbf{M}}-2\mathbf{C}$ 的 $\mathbf{C}$ 须取 Christoffel 形式]
若去掉重力补偿（$\boldsymbol{\tau}=-\mathbf{K}_p\mathbf{e}-\mathbf{K}_d\dot{\mathbf{e}}$），闭环右端变为 $\mathbf{g}(\mathbf{q}_d)-\mathbf{g}(\mathbf{q})\neq\mathbf{0}$（除非 $\mathbf{q}=\mathbf{q}_d$），平衡点漂成\emph{非零稳态误差}——这正是重力补偿在关节控制中是标配的原因。另外整套证明的命门是 $\dot{\mathbf{e}}^\top(\tfrac12\dot{\mathbf{M}}-\mathbf{C})\dot{\mathbf{e}}=0$，它依赖 $\mathbf{C}$ 取\emph{Christoffel 符号}参数化（$\mathbf{C}$ 的选择不唯一）；随意选 $\mathbf{C}$ 可能破坏反对称、使证明失效。再者若把 $\mathbf{K}_d$ 换成干摩擦 $\mathbf{K}_d\operatorname{sgn}(\dot{\mathbf{e}})$，$\mathbf{f}$ 不连续，LaSalle 同 \cref{sec:lyap-lasalle} 那样不再直接适用。
\end{pitfall}

\begin{insight}[与无源性/阻抗控制的接口]
$V=$ 总能量、$\dot V=-$ 阻尼耗散这一论证本质是\emph{无源性}：闭环把存储能量 $V$ 与耗散 $\dot{\mathbf{e}}^\top\mathbf{K}_d\dot{\mathbf{e}}$ 对接。这条能量论证在操作空间动力学与笛卡尔阻抗控制（\cref{ch:operational-space,ch:force-wbc}）中升级为“无源性塑形”；若内环用计算力矩（\cref{sec:ctrl-robot-ctc}）把机器人线性化为双积分器，则 \cref{ex:lyap-backstepping} 的反步/外环-内环设计可直接套用。
\end{insight}

% ════════════════════════════════════════════════════════════════
\section{吸引域估计与 Zubov 方程}
\label{sec:lyap-roa}

\paragraph{动机：吸引域有多大？}
即便证了渐近稳定，工程上关键的下一问是：\emph{吸引域 $\mathcal{R}=\{\mathbf{x}_0:\mathbf{x}(t;\mathbf{x}_0)\to\mathbf{0}\}$ 有多大？}机器人受大扰动飞出吸引域，控制器就失效。精确算 $\mathcal{R}$ 通常不可能，但 Lyapunov 给\emph{内估计}。

\paragraph{sublevel set 内估计与 SOS-ROA。}
若 $\Omega_c=\{V\le c\}$ 紧且 $\subset\{\dot V<0\}$，则 $\Omega_c\subset\mathcal{R}$；取最大的 $c$ 使 $\Omega_c$ 仍含于 $\{\dot V<0\}$ 即得最大内估计。对多项式系统可表为 SOS 优化
\begin{equation}
  \max_{\rho,V}\ \rho\quad\text{s.t.}\quad V-\epsilon\|\mathbf{x}\|^2\ \text{SOS},\quad -\dot V-\lambda(\rho-V)\ \text{SOS}\ (\lambda\ \text{SOS 乘子}),
  \label{eq:lyap-sos-roa}
\end{equation}
其中 S-procedure 乘子 $\lambda$ 把“在 $\{V\le\rho\}$ 内 $\dot V<0$”编码为 SOS 约束。

\begin{theorem}[Zubov 方程；Zubov 1964]
\label{thm:lyap-zubov}
存在 $V:\mathcal{R}\to[0,1)$ 与正定 $\phi$ 满足一阶 PDE
\begin{equation}
  \nabla V(\mathbf{x})\cdot\mathbf{f}(\mathbf{x})=-\phi(\mathbf{x})\,(1-V(\mathbf{x}))\,\sqrt{1+\|\mathbf{f}(\mathbf{x})\|^2},
  \label{eq:lyap-zubov}
\end{equation}
使吸引域被精确刻画为 $\mathcal{R}=\{V<1\}$，$\partial\mathcal{R}=\{V=1\}$\cite{zubov1964methods}。
\end{theorem}

\begin{insight}[Zubov $\leftrightarrow$ HJB；ROA $\leftrightarrow$ CBF]
Zubov 方程是一阶非线性 PDE，通常数值求解；$\sqrt{1+\|\mathbf{f}\|^2}$ 是标准归一化变体。它可视为\emph{零控制 HJB 的黏性解形式}——与 \cref{ch:dp-hjb} 的 HJB 黏性解、\cref{ch:hj-reachability} 的可达集互为镜像（值函数水平集即吸引域/可达域边界）。另一方面，ROA 的 sublevel set 内估计与 \emph{CBF 安全证书}是稳定性对偶：Lyapunov/ISS 保证“收敛进一个集”，CBF 保证“永不出一个集”，二者都用 sublevel/superlevel set + SOS 证书（CBF-QP 综合在 \cref{ch:clf-cbf}，多机/可达侧在 P3 规划部已展开，本章不重复）。
\end{insight}

% ════════════════════════════════════════════════════════════════
\section{非自治系统、UAS 与 Barbalat 引理}
\label{sec:lyap-nonauto}

\paragraph{动机：时变系统要额外小心。}
机器人大量场景\emph{非自治}：时变轨迹跟踪（$\mathbf{q}_d(t)$ 使闭环显含 $t$）、周期切换步态、衰减增益的自适应控制。对 $\dot{\mathbf{x}}=\mathbf{f}(\mathbf{x},t)$，经典定理须加\emph{一致性}要求，且 LaSalle 失效——这是 \cref{ch:control}（\cref{thm:ctrl-indirect}）仅一句话提及、本章系统补全之处。

\begin{definition}[一致渐近稳定 UAS]
\label{def:lyap-uas}
$\dot{\mathbf{x}}=\mathbf{f}(\mathbf{x},t)$ 原点 UAS，若 $\exists\beta\in\mathcal{KL}$ 使对所有 $t_0\ge0$：$\|\mathbf{x}(t)\|\le\beta(\|\mathbf{x}(t_0)\|,\,t-t_0),\ \forall t\ge t_0$。“一致”指 $\beta$ \emph{不依赖} $t_0$——无论何时启动，衰减行为都一样。
\end{definition}

\begin{theorem}[非自治 Lyapunov 定理；Khalil Thm 4.9]
\label{thm:lyap-nonauto}
若 $V(\mathbf{x},t)$ 满足\emph{一致}上下夹 $\alpha_1(\|\mathbf{x}\|)\le V(\mathbf{x},t)\le\alpha_2(\|\mathbf{x}\|)\ (\alpha_1,\alpha_2\in\mathcal{K}_\infty)$ 且 $\dot V=\frac{\partial V}{\partial t}+\nabla_{\mathbf{x}}V\cdot\mathbf{f}\le-W(\mathbf{x})$（$W$ 正定），则原点 UAS。
\end{theorem}

\noindent 关键区别在“一致夹对所有 $t$ 成立”：自治时 $V(\mathbf{x})$ 不含 $t$，夹自动一致；非自治须显式验证。若 $V$ 随时间太剧烈（如 $V=e^tx^2$ 上界不一致，或 $V=e^{-t}x^2$ 下界不一致），定理失效。

\begin{theorem}[Barbalat 引理；Barbalat 1959]
\label{thm:lyap-barbalat}
若 $f:[0,\infty)\to\mathbb{R}$ 一致连续且 $\lim_{t\to\infty}\int_0^t f(\tau)\,\mathrm d\tau$ 存在有限，则 $\lim_{t\to\infty}f(t)=0$\cite{barbalat1959systemes}。
\end{theorem}

\begin{insight}[Barbalat 是非自治的“LaSalle 替代”]
LaSalle（\cref{thm:lyap-lasalle-full}）\emph{只对自治系统成立}——其证明依赖 $\omega$-极限集的不变性，时变系统没有这个结构。非自治时用 Barbalat：若能证 $\dot V\le0$ 且 $\ddot V$ 有界（保证 $\dot V$ 一致连续），则 $\dot V(t)\to0$；若 $\dot V=0$ 蕴含某误差变量趋零即得收敛。典型应用：自适应系统 $\dot x=-x+\theta y,\ \dot\theta=-xy$，取 $V=\tfrac12(x^2+\theta^2)$ 得 $\dot V=-x^2\le0$，由 Barbalat 推 $x(t)\to0$（而参数 $\theta$ 未必收敛——这正是自适应控制“跟踪误差趋零但参数未必辨识到真值”的根源）。务必记住：Barbalat 给的是 $\dot V\to0$（沿时间），\emph{不}是“轨迹趋向不变集”，结论弱于 LaSalle，但对非自治是唯一选择。
\end{insight}

% ════════════════════════════════════════════════════════════════
\section{展望：超越平衡点}
\label{sec:lyap-outlook}

\paragraph{从平衡点到轨迹。}
经典 Lyapunov 以\emph{平衡点}为中心，但很多机器人任务的目标是“跟踪一条轨迹”而非“趋向一个点”（步态、末端轨迹跟踪）。\emph{Contraction Theory}（Lohmiller–Slotine 1998\cite{lohmiller1998contraction}）把 Lyapunov 推广为\emph{增量/微分}稳定性：系统 $\dot{\mathbf{x}}=\mathbf{f}(\mathbf{x},t)$ 称收缩，若存在黎曼度量 $\mathbf{M}(\mathbf{x},t)\succ0$ 使广义雅可比满足 $\mathbf{F}+\mathbf{F}^\top\preceq-2\lambda\mathbf{I}\ (\lambda>0)$，则\emph{任意两条轨迹}间的测地距离以 $e^{-\lambda t}$ 收缩——\emph{不需要知道平衡点在哪里}。直观地，Lyapunov 看 $V(\mathbf{x})$（关于状态），Contraction 看 $\delta\mathbf{x}^\top\mathbf{M}\delta\mathbf{x}$（关于两轨迹偏差）；有平衡点时 Contraction 退化为指数稳定。它天然适合轨迹跟踪、同步、观测器（任意初始估计指数收敛），并能处理极限环——综述见 Tsukamoto–Chung–Slotine 2021\cite{tsukamoto2021contraction}。

\paragraph{另两条升级方向（指针）。}
\emph{切换/混合系统}：即使每个子系统单独稳定，快速切换也可能发散；稳定性须靠\emph{共同 Lyapunov 函数}（任意切换下 GAS，对应 LMI $\mathbf{A}_i^\top\mathbf{P}+\mathbf{P}\mathbf{A}_i\prec0\ \forall i$ 可行）或\emph{多 Lyapunov 函数 + dwell-time}（切换间隔够大）——本书腿足部 P7 在步态切换语境展开。\emph{随机 Lyapunov}：确定性 $\dot V=\nabla V\cdot\mathbf{f}$ 升级为 Itô 生成元 $\mathcal{L}V=\nabla V\cdot\mathbf{b}+\tfrac12\operatorname{tr}(\boldsymbol{\sigma}\boldsymbol{\sigma}^\top\nabla^2 V)$，多出的二阶项 $\tfrac12\operatorname{tr}(\boldsymbol{\sigma}\boldsymbol{\sigma}^\top\nabla^2 V)$ 来自布朗运动二次变差 $(\mathrm dW)^2=\mathrm dt$，是随机机会约束/随机 MPC（P3 规划部）的判稳基础。这两者与本章主线（连续系统平衡点稳定性）耦合较弱，故仅作指针。

\paragraph{启下。}
本章把稳定性从“会用”做到了“会证、会扩”：比较函数立语言、直接法/LaSalle 证到底、指数稳定接 Riccati、逆定理给完备性、ISS/小增益/级联/反步建鲁棒与设计、PD 旗舰落地、ROA/非自治补全工程拼图。它直接喂给三个兄弟章：ISS-Lyapunov + 比较函数是 \cref{ch:clf-cbf}（CLF/CBF 安全控制）的前置；ISS 小增益的频域线性版在 \cref{ch:robust-hinf}（鲁棒/$H_\infty$）；终端代价=闭环 Lyapunov 函数支撑 \cref{ch:mpc-advanced}（MPC 深化）。带着“构造证书”的眼光，我们便能进入后续各章的控制综合。

% ════════════════════════════════════════════════════════════════
\section*{本章小结}
\addcontentsline{toc}{section}{本章小结}

\begin{table}[H]\centering\small
\caption{Lyapunov 稳定性理论：知识点速查}
\label{tab:lyap-summary}
\begin{tabular}{lll}
\toprule
知识点 & 核心条件/结论 & 关键应用 \\
\midrule
比较函数 & $\mathcal{K}/\mathcal{K}_\infty/\mathcal{KL}$ 三类 & 稳定性统一语言、ISS 货币 \\
直接法（完整证） & 正定 $V$ + 负（半）定 $\dot V$；$\varepsilon$-$\delta$+$\omega$-极限 & 一切稳定性证明母本 \\
LaSalle（完整证） & 紧正不变集 + 最大不变集 $M$ & 机械系统、PD 控制 \\
指数稳定 & 二次夹 $\Rightarrow k,\lambda$；Lyap 方程积分解 & 线性/局部、收敛率 \\
Chetaev & 锥内 $V>0,\dot V>0$ & 不稳定判定、倒立摆 \\
构造法 & 能量/二次型/Krasovskii/变梯度/SOS/Neural & 实际系统分析 \\
逆定理 & AS $\Leftrightarrow$ $\exists C^\infty V$ & 哲学完备性、保证搜索目标存在 \\
ISS & $\|\mathbf{x}\|\le\beta+\gamma(\|\mathbf{u}\|_\infty)$；ISS-Lyap & 带扰动鲁棒性 \\
小增益 & $\gamma_1\circ\gamma_2<\mathrm{id}$ & 互联/级联/分布式 \\
级联 & 上层 GAS + 下层 GAS + ISS & MPC+底层控制 \\
反步法 & 虚拟控制 $\phi$ + 增广 $V_2$ & 非线性控制器设计 \\
PD+重力补偿 & $\dot V=-\dot{\mathbf{e}}^\top\mathbf{K}_d\dot{\mathbf{e}}$ + LaSalle & 机器人调节（GAS） \\
ROA/Zubov & sublevel set/SOS/Zubov PDE & 安全验证、接 HJB \\
非自治/Barbalat & 一致夹 + UAS；Barbalat 替代 LaSalle & 轨迹跟踪、自适应 \\
\bottomrule
\end{tabular}
\end{table}

\section*{常见误解}
\addcontentsline{toc}{section}{常见误解}
\begin{enumerate}
  \item \textbf{“$\dot V\le0$ 就稳定”}：必须 $V$ \emph{正定}（非半正定），否则 $V=0$ 的非零点令论证崩溃。
  \item \textbf{“找到 $V$ 即 GAS”}：仅得局部；GAS 须 $V$ 径向无界且 $\dot V<0$ 在\emph{整个} $\mathbb{R}^n$ 成立。
  \item \textbf{“半负定 $\dot V$ 直接判渐近”}：须走 LaSalle，证 $\{\dot V=0\}$ 内最大不变集只含原点。
  \item \textbf{“线性化稳定即非线性稳定”}：间接法仅局部；临界（零实部）失效，须直接法。
  \item \textbf{“找不到 $V$ 即不稳定”}：逆定理反驳——稳定系统必有 $V$，找不到是技巧/算力问题。
  \item \textbf{“PB 给极限环唯一性”}：PB 只给\emph{存在性}，唯一性须 Liénard/Levinson–Smith。
  \item \textbf{“LaSalle 用于非自治”}：LaSalle 只对自治；非自治用 Barbalat。
  \item \textbf{“Krasovskii 直接给 GAS”}：$\mathbf{J}+\mathbf{J}^\top\preceq-\alpha\mathbf{I}$ 只给局部 AS；GAS 须补 $\|\mathbf{f}\|$ 径向无界。
\end{enumerate}

\section*{延伸阅读}
\addcontentsline{toc}{section}{延伸阅读}
\begin{itemize}
  \item \textbf{Khalil, \emph{Nonlinear Systems}, 3ed}\cite{khalil2002nonlinear}：本章主参考，定理 4.1/4.2/4.4/4.7/4.9/4.10、Lemma 4.4 的证明最完整。
  \item \textbf{Sontag, \emph{Mathematical Control Theory}, 2ed}\cite{sontag1998mathematical}：ISS 创始人视角，iISS 与等价刻画（免费 PDF）。
  \item \textbf{Slotine \& Li, \emph{Applied Nonlinear Control}}\cite{slotine1991applied}：工程导向，Barbalat 与反步法的简洁处理。
  \item \textbf{Krstić, Kanellakopoulos \& Kokotović, \emph{Nonlinear and Adaptive Control Design}}\cite{krstic1995nonlinear}：反步法经典专著。
  \item \textbf{Tsukamoto, Chung \& Slotine, \emph{Annual Reviews in Control} 2021}\cite{tsukamoto2021contraction}：Contraction Theory 最新综述（接 \cref{sec:lyap-outlook}）。
\end{itemize}
